% IN THE NAME OF GOD
% LaTeX Problem Set Template by Sachin Padmanabhan edited by Ali Heydari
% I created this when I was a freshman in CS 103,
% and I continue to use it to this day.
%
% Hope you enjoy!
%
% There may be problems with this template.
% If so, feel free to contact me.
%

\documentclass[a4paper]{article}
\usepackage{amsmath}
\usepackage{amssymb}
\usepackage{amsthm}
\usepackage{amssymb}
\usepackage{mathdots}
\usepackage[pdftex]{graphicx}
\usepackage{fancyhdr}
\usepackage[margin=1in]{geometry}
\usepackage{multicol}
\usepackage{bm}
\usepackage{listings}
\PassOptionsToPackage{usenames,dvipsnames}{color}  %% Allow color names
\usepackage{pdfpages}
\usepackage{algpseudocode}
\usepackage{tikz}
\usepackage{enumitem}
\usepackage[T1]{fontenc}
\usepackage{inconsolata}
\usepackage{framed}
\usepackage{wasysym}
\usepackage[thinlines]{easytable}
\usepackage{hyperref}
\usepackage{dsfont}


\usepackage{wrapfig}
\setlength{\intextsep}{0pt}
\setlength{\columnsep}{0pt}
\usepackage{subcaption}
\usepackage{graphicx}
\graphicspath{ {images/} }

\hypersetup{
    colorlinks=true,
    linkcolor=blue,
    filecolor=magenta,
    urlcolor=blue,
}

\title{
\textsc{
\center
\includegraphics[scale=0.5]{IUST_logo_color} \\
Iran university of science and technology} \\ [25pt] % Your university, school and/or department name(s)
Department of Computer Engineering \\
Discrete Mathematics\\Problem Set \#5 \\
}
\author{Ali Heydari}
\date{\today}

\lhead{Ali Heydari}
\chead{Problem Set \#5}
\rhead{\today}
\lfoot{}
\cfoot{Discrete mathematics --- Spring 2018}
\rfoot{\thepage}

\newcommand{\abs}[1]{\lvert #1 \rvert}
\newcommand{\absfit}[1]{\left\lvert #1 \right\rvert}
\newcommand{\norm}[1]{\left\lVert #1 \right\rVert}
\newcommand{\eval}[3]{\left[#1\right]_{#2}^{#3}}
\renewcommand{\(}{\left(}
\renewcommand{\)}{\right)}
\newcommand{\floor}[1]{\left\lfloor#1\right\rfloor}
\newcommand{\ceil}[1]{\left\lceil#1\right\rceil}
\newcommand{\pd}[1]{\frac{\partial}{\partial #1}}
\newcommand{\inner}[1]{\langle#1\rangle}
\newcommand{\cond}{\bigg|}
\newcommand{\rank}[1]{\mathbf{rank}(#1)}
\newcommand{\range}[1]{\mathbf{range}(#1)}
\newcommand{\nullsp}[1]{\mathbf{null}(#1)}
\newcommand{\repr}[1]{\left\langle#1\right\rangle}

\DeclareMathOperator{\Var}{Var}
\DeclareMathOperator{\tr}{tr}
\DeclareMathOperator{\Tr}{\mathbf{Tr}}
\DeclareMathOperator{\diag}{\mathbf{diag}}
\DeclareMathOperator{\dist}{\mathbf{dist}}
\DeclareMathOperator{\prob}{\mathbf{prob}}
\DeclareMathOperator{\dom}{\mathbf{dom}}
\DeclareMathOperator{\E}{\mathbf{E}}
\DeclareMathOperator{\R}{\mathbb{R}}
\DeclareMathOperator{\var}{\mathbf{var}}
\DeclareMathOperator{\quartile}{\mathbf{quartile}}
\DeclareMathOperator{\conv}{\mathbf{conv}}
\DeclareMathOperator{\VC}{VC}
\DeclareMathOperator*{\argmax}{arg\,max}
\DeclareMathOperator*{\argmin}{arg\,min}
\DeclareMathOperator{\Ber}{Bernoulli}
\DeclareMathOperator{\NP}{\mathbf{NP}}
\DeclareMathOperator{\coNP}{\mathbf{coNP}}
\DeclareMathOperator{\TIME}{\mathsf{TIME}}
\DeclareMathOperator{\polytime}{\mathbf{P}}
\DeclareMathOperator{\PH}{\mathbf{PH}}
\DeclareMathOperator{\SIZE}{\mathbf{SIZE}}
\DeclareMathOperator{\ATIME}{\mathbf{ATIME}}
\DeclareMathOperator{\SPACE}{\mathbf{SPACE}}
\DeclareMathOperator{\ASPACE}{\mathbf{ASPACE}}
\DeclareMathOperator{\NSPACE}{\mathbf{NSPACE}}
\DeclareMathOperator{\Z}{\mathbb{Z}}
\DeclareMathOperator{\N}{\mathbb{N}}
\DeclareMathOperator{\EXP}{\mathbf{EXP}}
\DeclareMathOperator{\NEXP}{\mathbf{NEXP}}
\DeclareMathOperator{\NTIME}{\mathbf{NTIME}}
\DeclareMathOperator{\DTIME}{\mathbf{DTIME}}
\DeclareMathOperator{\poly}{poly}
\DeclareMathOperator{\BPP}{\mathbf{BPP}}
\DeclareMathOperator{\ZPP}{\mathbf{ZPP}}
\DeclareMathOperator{\RP}{\mathbf{RP}}
\DeclareMathOperator{\coRP}{\mathbf{coRP}}
\DeclareMathOperator{\BPL}{\mathbf{BPL}}
\DeclareMathOperator{\IP}{\mathbf{IP}}
\DeclareMathOperator{\PSPACE}{\mathbf{PSPACE}}
\DeclareMathOperator{\NPSPACE}{\mathbf{NPSPACE}}
\DeclareMathOperator{\SAT}{\mathsf{SAT}}
\DeclareMathOperator{\NL}{\mathbf{NL}}
\DeclareMathOperator{\PCP}{\mathbf{PCP}}
\DeclareMathOperator{\PP}{\mathbf{PP}}
\DeclareMathOperator{\cost}{cost}
\let\Pr\relax
\DeclareMathOperator*{\Pr}{\mathbf{Pr}}

\definecolor{shadecolor}{gray}{0.95}

\theoremstyle{plain}
\newtheorem*{lem}{Lemma}

\theoremstyle{plain}
\newtheorem*{claim}{Claim}

\theoremstyle{definition}
\newtheorem*{answer}{Answer}

\newtheorem{theorem}{Theorem}[section]
\newtheorem*{thm}{Theorem}
\newtheorem{corollary}{Corollary}[theorem]
\newtheorem{lemma}[theorem]{Lemma}

\renewcommand{\headrulewidth}{0.4pt}
\renewcommand{\footrulewidth}{0.4pt}

\setlength{\parindent}{0pt}

\pagestyle{fancy}

\renewcommand{\thefootnote}{\fnsymbol{footnote}}

\begin{document}
\maketitle

\section{Connected graph}
Prove the theorem blow.
\begin{theorem}
A graph is connected if and only if for every partition of it's vertices into two nonempty sets, there
is an edge with endpoints in both sets.
\end{theorem}
\begin{shaded}
\begin{answer}
ANSWER HERE.
\end{answer}
\end{shaded}

\section{$n$-vertex graph}
Prove the theorem blow.
\begin{theorem}
Every $n$-vertex graph with at least $n$ edges contains a cycle.
\end{theorem}
\begin{shaded}
\begin{answer}
ANSWER HERE.
\end{answer}
\end{shaded}
\section{Connected simple $n$-vertex graph}
Prove the theorem blow.
\begin{theorem}
If $l, m, n$ are nonnegative integers with $l + m = n \geq 1$, then there exists a connected simple $n$-vertex
graph with $l$ vertices of even degree and $m$ vertices of odd degree if and only if $m$ is even,
except for $(l, m, n) = (2, 0, 2)$.
\end{theorem}
\begin{shaded}
\begin{answer}
ANSWER HERE.
\end{answer}
\end{shaded}



\section{$k\_cube$ graph}
what is $k\_cube(Qk)$ graph?

Count number of it's Vertices and Edges and prove that it is bipartite.
\begin{shaded}
\begin{answer}
\begin{proof}
PROOF HERE.
\end{proof}
\end{answer}
\end{shaded}

\section{Diameter of graph}
diameter of $G$ is length of the longest path between two vertices in it.
Show that if $G$'s diameter is greater than 3 its complement's diameter would be less than 3.
\begin{shaded}
\begin{answer}
\begin{proof}
PROOF HERE.
\end{proof}
\end{answer}
\end{shaded}

\section{About Graph}
Necessary and sufficient conditions for a list $d$ to be graphic when $d$ consists of $k$ copies of $a$ and $n − k$ copies of $b$, with $a \geq b \geq 0$. Since the degree sum must be even, the quantity $k$ $a + (n − k)b$ must be even. In addition, the inequality $k$ $a \leq k(k − 1) + (n − k)$ $min\{k, b\}$ must hold, since each vertex with degree $b$ has at most $min\{k, b\}$ incident edges whose other endpoint has degree $a$. We construct graphs with the desired degree sequence when these conditions hold. Note that the inequality implies $a \leq n − 1$.
\begin{shaded}
\begin{answer}
\begin{proof}
PROOF HERE.
\end{proof}
\end{answer}
\end{shaded}

\section{Graph orientation}
Prove the theorem blow.
\begin{theorem}
Every graph $G$ has an orientation such that $|d +(v) - d-(v)| \leq 1$ for all $v$.
\end{theorem}
\begin{shaded}
\begin{answer}
ANSWER HERE.
\end{answer}
\end{shaded}


\section{Strong orientation of a graph}
A strong orientation of a graph that has an odd cycle also has an odd $(directed)$ cycle. Suppose that $D$ is a strong orientation of a graph $G$ that has an odd cycle $v1, \dots , v2k+1$. Since $D$ is strongly connected, for each $i$ there is a $vi, vi+1-path$ in $D$. If for some $i$ every such path has even length, then the edge between $vi$ and $vi+1$ points from $vi+1$ to $vi$, since the other orientation would be a $vi, vi+1-path$ of length 1 $(odd)$. In this case, we have an odd cycle through $vi$ and $vi+1$. Otherwise, we have a path of odd length from each $vi$ to $vi+1$. Combining these gives a closed trail of odd length In a digraph as well as in a graph (by the same proof), a closed odd trail contains the edges of an odd cycle
\begin{shaded}
\begin{answer}
ANSWER HERE.
\end{answer}
\end{shaded}

\section{The Tournaments}
Tournaments with all players kings.

 If $n$ is odd, then there is an tournament with $n$ vertices such that every player is a king.
\begin{shaded}
\begin{answer}
ANSWER HERE.
\end{answer}
\end{shaded}


\section{About Tree}
For a tree $T$ with vertex degrees in ${1, k}$, the possible values of $n(T)$ are the positive integers that are 2 more than a multiple of $k − 1$.
\begin{shaded}
\begin{answer}
ANSWER HERE.
\end{answer}
\end{shaded}



\section{More about Tree}
Prove the theorems blow.
\begin{enumerate}[label=\Roman*]
  \item {
  \begin{theorem}
A tree has exactly one center or has two adjacent centers.
\end{theorem}
\begin{shaded}
\begin{answer}
ANSWER HERE.
\end{answer}
\end{shaded}
}

  \item {
  \begin{theorem}
A tree has exactly one center if and only if its diameter is twice its radius.
\end{theorem}
\begin{shaded}
\begin{answer}
ANSWER HERE.
\end{answer}
\end{shaded}
}
\end{enumerate}

\section{Cutting edge}
Prove the theorem blow.
\begin{theorem}
$e = xy$ is a cutting edge if and only if
 $
\begin{cases}
  G-e \quad \text{does not have ha} \quad x \rightarrow y \quad \text{path}. \\
\text{it does not belong to any cycle.}
\end{cases}
$
\end{theorem}
\begin{shaded}
\begin{answer}
ANSWER HERE.
\end{answer}
\end{shaded}

\section{$5-cycles$ in peterson graph}
how many $5-cycles$ does \textbf{peterson} graph have?
\begin{shaded}
\begin{answer}
ANSWER HERE.
\end{answer}
\end{shaded}

\section{Certain bridge club}
A certain bridge club has a special rule to the effect that four members may play together only if no two of them have previously partnered one another. At one meeting fourteen members, each of whom has previously partnered five others, turn up. Three games are played, and then proceedings come to a halt because of the club rule. Just as the members are preparing to leave, a new member, unknown to any of them, arrives. Show that at least one more game can now be played.
\begin{shaded}
\begin{answer}
ANSWER HERE.
\end{answer}
\end{shaded}

\end{document} 